\newpage
\subsection{Messen der Eingangsruheströme}\label{sec:2-1}
\subsubsection{Aufgabenstellung}
Es sei mit einer geeigneten Schaltung die Eingangsruheströme ($I_{Bn}$ und $I_{Bp}$) des OP-A3 ($\mu$A741 Biploar) auf dem Messbrett mittels geeigneter Schalung ermittelt werden.

\subsubsection{Versuchsaufbau}
Zur Messung von $I_{Bn}$ wird die Schaltung aus Abbildung~\ref{fig:ibn} verwendet. An $U_a$ wird das Osziloskop (siehe Geräteliste~\ref{sec:geraeteliste}) geschlossen. Sobald der Schalter $S_1$ geöffnet wird lädt sich $C$ mit $I_{Bn}$ linear auf. Der Schalter ist notwendig um den Kondensator vor der Messung zu entladen.

Ähnlich ist es mit $I_{Bp}$. Hier wird die Schalung aus Abbildung~\ref{fig:ibp} verwendet. Auch hier lädt sich der Kondensator $C$ nach dem öffnen des Schalters $S_1$ linear auf. Der OP wird als Impedanzwandler betrieben.

Alle Bauteile sind auf dem Messverstärkerbrett (siehe Geräteliste~\ref{sec:geraeteliste})

\begin{figure}[h]
	\begin{minipage}[position]{0.48\textwidth}
		\centering
		\begin{circuitikz}[european, scale=0.8, transform shape]
			% --- Op-amp A3 ---
			\node[op amp] (A3) at (0,0) {};
			\node at ($(A3.center)+(0.8,0.3)$) {\small $A_3$};

			% Supply arrows (+12V, -12V)
			\draw[->] ($(A3.center)+(0.2,0.4)$) -- +(0,0.5) node[above]{\small $+12$V};
			\draw[->] ($(A3.center)+(0.2,-0.4)$) -- +(0,-0.5) node[below]{\small $-12$V};

			% --- Output node ---
			\coordinate (out) at ($(A3.out)+(1.0,0)$);
			\draw (A3.out) -- (out);

			%capacitor and switch
			\coordinate (a-neg-mid) at ($(A3.-)+(-0.2,0)$);
			\coordinate (a-out-mid) at ($(A3.out)+(0.2,0)$);
			\coordinate (c-in) at ($(a-neg-mid)+(0,2.7)$);
			\coordinate (c-out) at ($(a-out-mid)+(0,3.2)$);
			\coordinate (switch-in) at ($(c-in)+(0,1.4)$);
			\coordinate (switch-out) at ($(c-out)+(0,1.4)$);
			\draw (c-in) to[C, l^={C}] (c-out);
			\draw (switch-in) to[ospst, l^={$S_1$}] (switch-out);
			\draw (switch-in) -- (a-neg-mid);
			\draw (switch-out) -- (a-out-mid);
			\draw[->] ($(c-in)+(+0.5,-0.7)$) -- ($(c-out)+(-0.5,-0.7)$)
			node[midway,below] {$U_C$};

			% V_+
			\node[ground] at (A3.+) {};

			% --- Non-inverting input: voltage source G1 with U_os ---
			\coordinate (srcbottom) at (-2,-1.7);
			\node[ground] at ($(srcbottom)+(0,0.5)$) {};
			\draw (srcbottom) to[current source, i_={$I_{Bn}$}] ++(0,2.2) coordinate (srcnode);
			\draw (srcnode) -- (A3.-);

			% Output 
			\node[ocirc] at (out) {};
			\draw[->] ($(out.south)+(0,-0.2)$) -- ++(0,-0.7)
			node[midway,right] {$U_a$}
			coordinate (mid);
			\node[ocirc] (C2) at ($(mid)+(0,-0.2)$) {};
			\draw (C2.south) -- ++(0,-0.0) node[ground] {};

		\end{circuitikz}
		\caption{Messschaltung für $I_{Bn}$}\label{fig:ibn}
	\end{minipage}\hfill
	\begin{minipage}[position]{0.48\textwidth}
		\centering
		\begin{circuitikz}[european, scale=0.8, transform shape]
			% --- Op-amp A3 ---
			\node[op amp, yscale=-1] (A3) at (0,0) {};
			\node at ($(A3.center)+(0.8,0.3)$) {\small $A_3$};

			% Supply arrows (+12V, -12V)
			\draw[->] ($(A3.center)+(0.2,0.4)$) -- +(0,0.5) node[above]{\small $+12$V};
			\draw[->] ($(A3.center)+(0.2,-0.4)$) -- +(0,-0.5) node[below]{\small $-12$V};

			% --- Output node ---
			\coordinate (out) at ($(A3.out)+(1.0,0)$);
			\draw (A3.out) -- (out);

			%capacitor and switch
			\coordinate (a-pos-mid) at ($(A3.+)+(-1.0,0)$);
			\coordinate (a-pos-far-left) at ($(a-pos-mid)+(-2,0)$);
			\coordinate (c-in) at ($(a-pos-far-left)+(0,0)$);
			\coordinate (c-out) at ($(a-pos-mid)+(0,0)$);
			\coordinate (switch-in) at ($(c-in)+(0,1.4)$);
			\coordinate (switch-out) at ($(c-out)+(0,1.4)$);
			\draw (c-in) to[C, l^={C}] (c-out);
			\draw (switch-in) to[ospst, l^={$S_1$}] (switch-out);
			\draw (switch-in) -- (a-pos-far-left);
			\draw (switch-out) -- (a-pos-mid);
			\draw (A3.+) -- (a-pos-mid);
			\node[ground] at ($(a-pos-far-left)+(0,0)$) {};
			\draw[<-] ($(c-in)+(+0.5,-0.7)$) -- ($(c-out)+(-0.5,-0.7)$)
			node[midway,below] {$U_C$};

			% V_-
			\coordinate (a-out-mid) at ($(A3.out)+(0.2,0)$);
			\draw (A3.-) -- ++(0,-1.3) -- ++(2.6,0) -- (a-out-mid);

			% ibp
			\coordinate (srcbottom) at ($(a-pos-mid)+(0,-2.3)$);
			\node[ground] at ($(srcbottom)+(0,0.5)$) {};
			\draw (srcbottom) to[current source, i_={$I_{Bp}$}] ++(0,2.2) coordinate (srcnode);
			\draw (srcnode) -- (a-pos-mid);

			% Output 
			\node[ocirc] at (out) {};
			\draw[->] ($(out.south)+(0,-0.2)$) -- ++(0,-0.7)
			node[midway,right] {$U_a$}
			coordinate (mid);
			\node[ocirc] (C2) at ($(mid)+(0,-0.2)$) {};
			\draw (C2.south) -- ++(0,-0.0) node[ground] {};

		\end{circuitikz}
		\caption{Messschaltung für $I_{Bp}$}\label{fig:ibp}
	\end{minipage}
\end{figure}

\subsubsection{Dimensionierung}
Laut Tabelle~\cite[S.~77, Abb.~3.13]{EMTPraktikum2025} ist ein $I_B = 80\,\mu\mathrm{A}$ zu erwarten. Zum messen ist ein Zeitraum $t=\SI{1}{\second}$ gut geeignet. Um die Verstärkung gut zu nützen soll der Kondensator nach der Zeit $t$ auf $\SI{10}{\volt}$ aufgeladen sein. Mit dem Kondensator von $\SI{10}{\nano\farad}$ ist man den $\SI{10}{\volt}$ wie man in Gleichung~\ref{eq:ib-kondensator} sehen kann sehr nahe.

Allgemein gilt für die Spannung am Kondensator wenn er an $t_0$ entladen war die folgende Formel.
\begin{gather}
	U_C = \frac{1}{C} \int_{0}^{t} i(\tau)\,\mathrm{d}\tau
\end{gather}

Mit konstanten Strom vereinfacht sich das
\begin{gather}
	U_C = \frac{I\cdot t}{C}, \quad \implies{} C = \frac{I \cdot t}{U_A} = \frac{\SI{80}{\nano\ampere} \cdot \SI{1}{\second}}{\SI{10}{\volt}} = \SI{8}{\nano\farad} \label{eq:ib-kondensator} \\[1.3em]
	\text{Gewählt wurde: }
	C = 10\,\mathrm{nF}
	\quad\Rightarrow
	U_C = 8\,\mathrm{V}.
\end{gather}

\subsubsection{Messung und Berechnung $I_{Bn}$}
In Abbildung~\ref{fig:ibn-messung} sieht man, dass $U_a$ innerhalb von $\SI{3}{\second}$ um $\SI{9}{\volt}$ angestiegen ist. Daraus wird $I_{Bn}$ berechnet.

\begin{gather}
	\Delta t = \SI{3}{\second}, \quad	\Delta U_A = \SI{9}{\volt}, \quad U_C=-U_A \\[1.5em]
	I_{Bn}=-C \, \frac{\Delta U_A}{\Delta t}
	= -\SI{10}{\nano\farad} \, \frac{\SI{9}{\volt}}{\SI{3}{\second}}
	= -\SI{30}{\nano\ampere}
\end{gather}

\begin{figure}[h]
	\begin{center}
		\includegraphics[width=0.65\textwidth]{IB-N-Messung.JPG}
	\end{center}
	\caption{Messung von $I_{Bn}$}\label{fig:ibn-messung}
\end{figure}

\subsubsection{Messung und Berechnung $I_{Bp}$}
In Abbildung~\ref{fig:ibp-messung} sieht man, dass $U_a$ innerhalb von $\SI{3}{\second}$ um $\SI{6.6}{\volt}$ angestiegen ist. Daraus wird $I_{Bp}$ berechnet.

\begin{gather}
	\Delta t = \SI{3}{\second}, \quad	\Delta U_A = \SI{6.6}{\volt}, \quad U_C=U_A \\[1.5em]
	I_{Bp}=C \, \frac{\Delta U_A}{\Delta t}
	= \SI{10}{\nano\farad} \, \frac{\SI{-6.6}{\volt}}{\SI{3}{\second}}
	= -\SI{22}{\nano\ampere}
\end{gather}

\begin{figure}[h]
	\begin{center}
		\includegraphics[width=0.65\textwidth]{IB-P-Messung.JPG}
	\end{center}
	\caption{Messung von $I_{Bp}$}\label{fig:ibp-messung}
\end{figure}

\subsubsection{Berechnung $I_B$ und $I_{\text{offset}}$}
\begin{gather}
	I_B = \frac{I_{Bp}+I_{Bn}}{2} = \frac{-\SI{22}{\nano\ampere}-\SI{30}{\nano\ampere}}{2} = \SI{-26}{\nano\ampere} \\
	I_{\text{offset}} = |I_{Bp}-I_{Bn}| = |-\SI{22}{\nano\ampere}+\SI{30}{\nano\ampere}| = \SI{4}{\nano\ampere}
\end{gather}

\subsubsection{Erkenntnisse}
Die \SI{-26}{\nano\ampere} von $I_B$ sind 3 mal weniger als die \SI{80}{\nano\ampere} aus der Tabelle~\cite[S.~77, Abb.~3.13]{EMTPraktikum2025}. Abweichungen zum Typwert sind plausibel durch Temperatur, Exemplarstreuung sowie Leckströme von Kondensator/Schalter und Messaufbau.
